\chapter{Bounded Total Derived Functors}
\label{ch:derived-functors}

Throughout this chapter, categories are abelian, and the functors used to
construct total derived functors are additive. An assertion that a category
has enough injectives or projectives applies to the source of the corresponding
construction. We use \(\mathrm K(\mathcal A)\) for its homotopy category,
\(D(\mathcal A)\) for its localization, and \(Q_{\mathcal A}\) for the
localization functor from \(\mathrm K(\mathcal A)\). The superscripts \(+\)
and \(-\) mean bounded below and bounded above, respectively. On the complex
side the bounds are on terms, and on the derived side they are on cohomology.
The comparison between these descriptions is Proposition~\ref{res:D06}.

\section{Replacements and computation of morphisms}

We first construct replacements of strictly bounded complexes. Here is the
projective construction in an abelian category with enough projectives.
For a complex \(X\) zero in degrees \(>N\), start with \(P^i=0\) for \(i>N\).
Suppose the part of \(P\to X\) above degree \(n\) has been constructed. Set
\[
  W^n=X^n\times_{Z^{n+1}X}Z^{n+1}P,
\]
using \(d_X^n\) and the already defined comparison on cycles. Choose an
epimorphism \(P^n\twoheadrightarrow W^n\) from a projective object. Its first
projection defines \(p^n\), and its second projection, followed by the cycle
inclusion, defines \(d_P^n\). The pullback identities give \(p^{n+1}d_P^n=d_X^np^n\)
and \(d_P^{n+1}d_P^n=0\).

Pulling back \(P^n\twoheadrightarrow W^n\) along \(Z^nX\to W^n\),
\(z\mapsto(z,0)\), gives \(Z^nP\twoheadrightarrow Z^nX\). Thus
\(W^n\to X^n\) is an epimorphism by the same assertion in degree \(n+1\),
so \(p^n\) is an epimorphism. Finally, the pullback used at degree \(n-1\)
shows that \(d_P^{n-1}\) surjects onto the inverse image of \(B^nX\) in
\(Z^nP\). Indeed \(X^{n-1}\twoheadrightarrow B^nX\) is an epimorphism,
and so is its pullback to that inverse image. Consequently the induced
map \(Z^nP/B^nP\to Z^nX/B^nX\) is both monic and epic. This proves it is an
isomorphism for every \(n\).

\begin{theorem}[Bounded-below injective replacement]\label{res:R01}
If \(\mathcal A\) has enough injectives, a bounded-below complex \(X\)
admits a termwise monic quasi-isomorphism \(X\to I\), with \(I\) termwise
injective and zero below the same lower bound. Every object of
\(D^+(\mathcal A)\) admits such an injective representative.
\end{theorem}
\begin{proof}
Apply the preceding construction in \(\mathcal A^{\mathrm{op}}\) to
the bounded-above complex with degree \(i\) term \((X^{-i})^{\mathrm{op}}\)
and degree \(i\) differential \((d_X^{-i-1})^{\mathrm{op}}\). This reversal
interchanges kernels and cokernels, so the comparison returns a
termwise monic quasi-isomorphism with the stated lower bound. Use
Proposition~\ref{res:D06} first if the input is an object of the derived
category. See \cite[Tag 013K]{stacks-derived}.
\end{proof}

\begin{theorem}[Bounded-above projective replacement]\label{res:R02}
If \(\mathcal A\) has enough projectives, a bounded-above complex \(X\)
admits a termwise epic quasi-isomorphism \(P\to X\), with \(P\) termwise
projective and zero above the same upper bound. Every object of
\(D^-(\mathcal A)\) admits such a projective representative.
\end{theorem}
\begin{proof}
The downward induction above constructs the epic comparison, and
Proposition~\ref{res:D06} supplies a bounded-above complex for a derived
object. The construction replaces a complex rather than only an object
in degree zero. Compare \cite[Tag 0646]{stacks-derived} and the ordinary
module resolutions of \cite[Chapter 10]{p103}.
\end{proof}

\begin{theorem}[Orthogonality of the bounded models]\label{res:R03}
A complex \(I\) is K-injective if every map to \(I\) from an acyclic
complex is null-homotopic. Dually \(P\) is K-projective if every map
from \(P\) to an acyclic complex is null-homotopic. A bounded-below
complex of injectives is K-injective, and a bounded-above complex of
projectives is K-projective.
\end{theorem}
\begin{proof}
For \(u:C\to I\) with \(C\) acyclic, construct
\(s^n:C^n\to I^{n-1}\) starting at the lower bound of \(I\). If the
homotopy identity has been established through degree \(n-1\), then
\(u^n-d_I^{n-1}s^n\) kills \(B^nC=Z^nC\) by that identity and the chain
map property of \(u\). It therefore factors through
\(\operatorname{im}d_C^n\subseteq C^{n+1}\). Extend this factorization
to \(s^{n+1}:C^{n+1}\to I^n\) by injectivity. This gives
\(u=d_Is+sd_C\) in every degree. Passing to the opposite category
and reversing cochain degrees gives the projective statement.
See \cite[Tags 013R, 0647]{stacks-derived}.
\end{proof}

\begin{theorem}[Morphisms computed on K-models]\label{res:R04}
For K-injective \(I\) and K-projective \(P\), the maps induced by
localization
\[
 \operatorname{Hom}_{\mathrm K}(X,I)\longrightarrow
     \operatorname{Hom}_{D}(QX,QI),\qquad
 \operatorname{Hom}_{\mathrm K}(P,X)\longrightarrow
     \operatorname{Hom}_{D}(QP,QX)
\]
are bijective and natural. A quasi-isomorphism between two K-injectives,
or between two K-projectives, is a homotopy equivalence.
\end{theorem}
\begin{proof}
For a quasi-isomorphism \(s:U\to X\), its cone is acyclic by
Proposition~\ref{res:K03}. Applying \(\operatorname{Hom}_{\mathrm K}(-,I)\)
to the cone triangle and using Proposition~\ref{res:K06} shows that
precomposition by \(s\) is a bijection. Thus a roof
\(X\xleftarrow{s}U\xrightarrow{u} I\) has a unique representative
\(v:X\to I\) in the homotopy category with \(vs=u\).
Refining two equal roofs leaves that representative unchanged, by the
same bijection. The right roof argument and
\(\operatorname{Hom}_{\mathrm K}(P,-)\) prove the dual assertion.
In each case an inverse to the image of a quasi-isomorphism is between
objects on which \(Q\) is fully faithful. It lifts to a two-sided inverse
in \(\mathrm K\), which is precisely the homotopy equivalence assertion.
This also supplies uniqueness up to homotopy of each resolution lift.
See \cite[Tags 05TG, 064B]{stacks-derived}.
\end{proof}

\section{Total functors and acyclic models}

Let \(F_K\) denote termwise application of \(F:\mathcal A\to\mathcal B\)
on either bounded homotopy category. It sends homotopies to homotopies.
Write \(U=Q_{\mathcal B}F_K\), so that \(U\) already takes values in the
target derived category. We use ordinary categorical composition order in
this chapter.

\begin{definition}[Total derived functor and its comparison]
\label{res:R05}
A total right derived functor comprises a functor \(RF\) on the
localized source and a natural transformation \(\eta:U\to RF Q_{\mathcal A}\).
For every functor \(H\) on the localized source, composition with \(\eta\)
identifies natural transformations \(RF\to H\) with transformations
\(U\to H Q_{\mathcal A}\). On the left, a comparison \(\epsilon:LF Q_{\mathcal A}\to U\)
identifies natural transformations \(H\to LF\) with transformations
\(H Q_{\mathcal A}\to U\).
\end{definition}

The comparisons characterize the functors up to unique compatible
isomorphism. In the right case, factor the second comparison through the
first and the first through the second. Both compositions induce the
original comparison, so both are identities by uniqueness. The left case
reverses the arrows. A natural transformation between the original
functors induces one between their right derived functors by the same
factorization. The left version uses the dual factorization.
See \cite[Tags 05SE, 05SV]{stacks-derived} and
\cite[Chapter 8]{yekutieli}.

\begin{theorem}[Total right derived functor on \(D^+\)]\label{res:R06}
With enough injectives in \(\mathcal A\), an additive \(F\) has an
exact total right derived functor \(RF:D^+(\mathcal A)\to D^+(\mathcal B)\).
It is computed by \(F(I)\) after replacement \(X\to I\) from
Theorem~\ref{res:R01}. The comparison on a model \(I\) is an isomorphism,
and natural transformations of additive functors are derived using
these same models.
\end{theorem}
\begin{proof}
The full homotopy subcategory of bounded-below complexes of injectives
is closed under shifts and cones, since their terms are finite biproducts
of injectives. Its composite with \(Q_{\mathcal A}\) is an exact
equivalence onto \(D^+(\mathcal A)\) by Theorems~\ref{res:R01}
and~\ref{res:R04}. Transport \(Q_{\mathcal B}F_K\) restricted to that
subcategory along an inverse equivalence. An additive functor
preserves the finite biproducts in the cone matrix and its signs, so the
transported functor is exact. Degreewise \(F(X\to I)\) gives \(\eta\),
which is natural because replacement lifts are unique up to homotopy
and \(F_K\) preserves homotopies.

To verify the universal property, a transformation \(U\to H Q_{\mathcal A}\)
can be restricted to the model subcategory. There \(\eta\) is an isomorphism,
so it gives the required map \(RF\to H\) via the equivalence. Naturality
on \(X\to I\) forces the original transformation on every complex,
which proves both the factorization and its uniqueness.
See \cite[Tags 05TI, 0159]{stacks-derived}.
\end{proof}

\begin{theorem}[Total left derived functor on \(D^-\)]\label{res:R07}
With enough projectives in \(\mathcal A\), an additive \(F\) has an
exact total left derived functor \(LF:D^-(\mathcal A)\to D^-(\mathcal B)\).
It is computed by \(F(P)\) after replacement \(P\to X\) from
Theorem~\ref{res:R02}.
\end{theorem}
\begin{proof}
Use the full homotopy subcategory of bounded-above complexes of
projectives. Shifts and cones preserve it. Theorems~\ref{res:R02}
and~\ref{res:R04} make its localization an exact equivalence onto
\(D^-(\mathcal A)\). Restrict \(Q_{\mathcal B}F_K\) and transport
across the equivalence. The comparison is induced by \(F(P\to X)\).
Given a transformation \(H Q_{\mathcal A}\to U\), restrict to projective
models, where this comparison is invertible, and transport back.
Naturality on \(P\to X\) shows it agrees with the given map for every \(X\).
This construction uses the same cochain shifts as the right one.
If expressed by duality, the reversal in Theorem~\ref{res:R01}
changes \([n]\) into \([-n]\) and interchanges the two bounds.
See \cite[Theorem 8.4.23]{yekutieli}.
\end{proof}

The unit and counit in these constructions commute with shifts:
the shift of a model is a model for the shifted complex, and
the resolution comparison comes from its shifted map. Uniqueness
of the comparison in the homotopy category makes this compatible
with using any choice of models.

\begin{theorem}[Computing with acyclic terms]\label{res:R08}
Call an object \(A\) right \(F\)-acyclic if the comparison
\(F(A)[0]\to RF(A[0])\) is an isomorphism, and define left
\(F\)-acyclicity using the reverse comparison to \(F(A)[0]\).
For a bounded-below complex \(C\) of right \(F\)-acyclic objects,
\(F(C)\to RF(C)\) is an isomorphism. For a bounded-above complex of left
\(F\)-acyclic objects, \(LF(C)\to F(C)\) is an isomorphism.
If \(F\) is left exact, right acyclicity reduces to the vanishing of the
positive right derived functors. For a right exact \(F\), the corresponding
left condition is the vanishing of the positive left derived functors.
\end{theorem}
\begin{proof}
In the right case begin with a finite complex. The degreewise split
sequence peeling off its bottom term gives a triangle by
Theorem~\ref{res:D04}. Application of \(F\) preserves this split sequence
degreewise. Naturality of the comparison gives a map of triangles,
and the long exact sequences show inductively that it is an isomorphism
on cohomology, starting with a shifted stalk.
Now fix a degree \(i\) and replace \(C\) by its brutal quotient truncation
in degrees \(\le i+1\). Its kernel has no terms below \(i+2\).
Both that kernel after \(F\) and its right derived image have
vanishing cohomology in degrees \(i,i+1\). For the derived image this
uses the bound-preserving replacement of Theorem~\ref{res:R01}.
The long exact sequences thus reduce comparison in degree \(i\)
to the already proved finite case.

For the left case, reverse the degrees and use the opposite categories.
The top-term sequence now peels off a top stalk, and the discarded
degrees are below those being tested. For a left exact \(F\), applying it
to the first two steps of an injective resolution identifies \(F(A)\)
with the zeroth cohomology of \(F(I)\). The replacement starts in degree
zero, so all negative cohomology vanishes. The right exact case uses the
first two projectives and cokernels. See
\cite[Tags 015E, 05TA]{stacks-derived}.
\end{proof}

\begin{theorem}[Classical derived functors and the delta-functor property]
\label{res:R09}
For a stalk \(A[0]\), the groups or objects
\(R^iF(A)=H^i(RF(A[0]))\) and \(L_iF(A)=H^{-i}(LF(A[0]))\)
are computed by ordinary injective and projective resolutions.
They vanish for \(i>0\) on injective and projective objects,
respectively. For left exact \(F\), \(R^0F=F\), and the right derived
functors form a universal cohomological delta-functor. For right exact
\(F\), \(L_0F=F\) and the left derived functors have the homological
analogue.
\end{theorem}
\begin{proof}
The replacements of a stalk are the ordinary resolutions, with a
projective resolution indexed in nonpositive cochain degrees. An
injective or projective stalk is itself an allowed model. The degree-zero
claims are those in Theorem~\ref{res:R08}. For a short exact sequence
of stalks use Theorem~\ref{res:D04} and apply the exact total functor.
The cohomology sequences give the connecting maps.

For universality on the right, for each object \(A\) embed it in an
injective \(I\) with quotient \(B\). The connecting map
\(R^{i-1}F(B)\to R^iF(A)\) is epic for \(i>0\), and its kernel is the
image from \(R^{i-1}F(I)\). Given a degree-zero transformation to
another delta-functor, construct its higher components using this
surjection and the corresponding connecting map of the target.
The previous degree of the transformation respects the image from \(I\),
so the construction factors through the surjection. For a morphism
\(A\to A'\), extend its composite \(A\to I'\) to \(I\to I'\).
Naturality of the two long exact sequences, and the surjection just
used, prove naturality and independence of the embedding. Compatibility
with connecting maps for an arbitrary short exact sequence follows by
mapping that sequence into \(0\to A\to I\to B\to0\), extending \(A\to I\)
over its middle object. Surjectivity also proves uniqueness.
Dually, given a degree-zero map from a homological delta-functor into
\(L_0F\), choose a projective epimorphism \(P\to A\) with kernel \(J\).
The inclusion \(L_iF(A)\to L_{i-1}F(J)\) for \(i>0\) forces the
higher maps by the same argument in the opposite category. See
\cite[Tags 015A, 05TE, 015B]{stacks-derived}.
\end{proof}

\begin{proposition}[Exact functors and localization]\label{res:R10}
An exact \(F\) descends termwise to an exact functor \(D(\mathcal A)\to D(\mathcal B)\)
commuting with cohomology. On the bounded derived categories on which
they exist, \(RF\) and \(LF\) identify with this functor.
\end{proposition}
\begin{proof}
Exactness identifies the kernels, images, and quotients defining
cohomology of \(F(X)\) with the \(F\)-images of those for \(X\).
It follows that \(F_K\) preserves quasi-isomorphisms, so the localization
property in Theorem~\ref{res:D03} gives the functor. Because
termwise application preserves cone matrices and shifts, the functor is
exact. Its comparison before localization is invertible, so it satisfies
both universal properties in Definition~\ref{res:R05} whenever
restricted to the relevant bounded category. See
\cite[Tag 015F]{stacks-derived}.
\end{proof}

\section{Ext, Tor, and comparison of compositions}

\begin{theorem}[Ext as shifted Hom]\label{res:R11}
For \(A,B\in\mathcal A\) the group
\(\operatorname{Hom}_{D(\mathcal A)}(A[0],B[n])\) agrees for \(n\ge0\)
with Ext computed by either of the available bounded resolutions,
covariantly in \(B\) and contravariantly in \(A\).
In degree zero it is \(\operatorname{Hom}_{\mathcal A}(A,B)\), and
it vanishes for \(n<0\).
\end{theorem}
\begin{proof}
The fully faithful heart embedding of Theorem~\ref{res:D07} gives
the zeroth formula. Its orthogonality gives the negative vanishing,
whether or not either kind of resolution is available.
For an injective resolution \(B\to I\), Theorem~\ref{res:R04}
computes the morphism group by homotopy classes of maps \(A[0]\to I[n]\).
Those maps are morphisms \(A\to I^n\) killed by the next differential,
modulo morphisms factoring through the previous differential via
\(A\to I^{n-1}\). The signs of \(I[n]\) do not change either subgroup.
This gives \(H^n\operatorname{Hom}(A,I^\bullet)\), the stated vanishing
and the zeroth identification. A projective resolution \(P\to A\) gives
the same argument for maps \(P\to B[n]\). In the usual cochain Hom
convention its differential on degree \(j\) is
\(d(u)=d_B u-(-1)^j u d_P\), so cycles are precisely maps to \(B[j]\);
boundaries are precisely homotopies to zero. If one writes the
resolution Hom differential as unsigned precomposition, multiply
its degree-\(j\) term by \((-1)^{j(j+1)/2}\) to identify the two
conventions. The Hom bijections are natural and lifts unique up to
homotopy, giving the bifunctor comparison. This identifies, in
particular, the resolution-based module Ext of \cite[Chapter 11]{p103}
with the derived model without identifying their representing complexes.
See \cite[Tags 06XQ--06XS]{stacks-derived}.
\end{proof}

\begin{theorem}[Right derived composition and acyclicity]\label{res:R12}
For left exact \(F:\mathcal A\to\mathcal B\) and \(G:\mathcal B\to\mathcal C\),
assume \(\mathcal A,\mathcal B\) have enough injectives. The canonical map
\(R(GF)\to RG\,RF\) on \(D^+(\mathcal A)\) is invertible if and only if
\(F(I)\) is right \(G\)-acyclic for every injective \(I\) of \(\mathcal A\).
\end{theorem}
\begin{proof}
Compose the comparison from \(G\) on the termwise complex \(F_K(X)\)
with \(RG\) of the comparison from \(F\) on \(X\). This is a
transformation \(Q_{\mathcal C}(GF)_K\to RG\,RF Q_{\mathcal A}\);
Definition~\ref{res:R05} induces the displayed map. Evaluate it on a
bounded-below injective model \(I^\bullet\) of \(X\). There its first
term is \(GF(I^\bullet)\), the intermediate comparison for \(F\)
is invertible, and it is exactly the comparison
\(G(F(I^\bullet))\to RG(F(I^\bullet))\).
The latter is invertible under the acyclicity condition by
Theorem~\ref{res:R08}. Conversely, evaluate an invertible comparison
at each injective stalk \(I[0]\). Its component is then precisely
\(G(F(I))[0]\to RG(F(I)[0])\). This is right \(G\)-acyclicity by
definition. See \cite[Tags 05T2, 015M]{stacks-derived}.
\end{proof}

\begin{proposition}[Ext boundaries and extension classes]\label{res:R13}
For a short exact sequence \(0\to A\to B\to C\to0\), write its derived
boundary as \(\delta:C\to A[1]\). The triangle map in the second Ext
variable is postcomposition with \(\delta[n]\); in the first it takes
\(u\) to \(u[1]\delta\). Under the injective Hom comparison of
Theorem~\ref{res:R11}, the ordinary second-variable boundary has the
displayed sign, and the ordinary first-variable boundary takes \(u\)
to \((-1)^{n+1}u[1]\delta\).
With extension classes given by the triangle convention of
Theorem~\ref{res:D04}, the assignment of \(\delta\) to a short exact
sequence identifies Baer sum with addition in \(\operatorname{Ext}^1(C,A)\).
\end{proposition}
\begin{proof}
The two representable long exact sequences of the derived triangle
give these maps, by the argument of Proposition~\ref{res:K06}.
For the second variable, regard \(A\to B\to C\) as a
three-term acyclic complex. Apply the Cartan--Eilenberg
construction of Theorem~\ref{res:G01}. Each successive
cokernel complex is acyclic by Lemma~\ref{res:G07}, so the
injective object for horizontal cohomology is always zero.
Each horizontal row is therefore a split short exact
sequence, and the three columns resolve \(A,B,C\).
This gives the injective horseshoe in this abelian category.
If its off-diagonal differential is
\(\alpha:I_C^j\to I_A^{j+1}\), the connecting class on a cocycle \(z\)
is \(\alpha z\), by taking its lift in the \(I_C\)-summand. The
map \(\delta\) is represented by \(\alpha:I_C^\bullet\to I_A^\bullet[1]\)
by Theorem~\ref{res:D04}. Thus postcomposition with \(\delta[n]\)
has the claimed sign.

For the first variable, use an injective resolution \(I^\bullet\) of
the fixed target, represent \(u\) by a cocycle \(u:A\to I^n\), and
extend it to \(w:B\to I^n\). Write \(d_Iw=v\pi\) for the quotient
\(\pi:B\to C\), so the usual Hom connecting map is represented by \(v\).
On the two-term cone \((A\xrightarrow{i}B)\) in degrees \(-1,0\),
the map \(u[1]\delta\) has components \((-u,0)\).
Adding the homotopy with degree-zero component \(w\) gives components
\((0,(-1)^{n+1}v\pi)\): the degree-zero target differential is that
of \(I[n+1]\), and the degree \(-1\) homotopy term is \(wi=u\).
This proves the comparison, with the sign shown in the statement.

Here is the extension classification in any abelian category.
By Theorem~\ref{res:D02} a morphism \(C[0]\to A[1]\) is represented
by a roof \(C[0]\leftarrow L^\bullet\to A[1]\) whose left arrow is
a quasi-isomorphism. Replace \(L\) first by \(\tau_{\le0}L\) and then by
its good quotient truncation \(\tau_{\ge-1}\). Both replacements
are quasi-isomorphisms and the arrows restrict and descend. For the
arrow to \(A[1]\), descent in degree \(-1\) follows since it kills
the boundaries there. The roof is now given by a short exact sequence
\(0\to L^{-1}\xrightarrow{d}L^0\to C\to0\) and a map \(a:L^{-1}\to A\).
Push this sequence out along \(-a\). The resulting cone roof maps
to \(A[1]\) by the negative projection, and hence represents the
given morphism. A homotopy change of \(a\) adds a map factoring
through \(d\). Pushing out along maps differing this way gives
equivalent extensions: in the biproduct presentation of the pushout,
add the corresponding multiple of the \(L^0\)-coordinate to the
\(A\)-coordinate. If two roofs agree, pass to a common refinement
as in Theorem~\ref{res:D02} and apply the same two truncations.
The left maps agree on cycles in degree zero, and the right maps
differ at most by the homotopy term just treated. The refined
extension therefore maps to both pushouts by maps inducing the
identity on \(A\) and \(C\). Conversely an equivalence of extensions
gives equal cone roofs. This proves the bijection.

Taking direct sums of the cone roofs, pulling back by the diagonal
of \(C\), and pushing out by the codiagonal of \(A\) sends two
boundaries to their sum. This is the Baer addition operation.
The projection convention of \cite[Tags 06XP, 06XU, 06XV]{stacks-derived}
uses the positive projection for the Yoneda class; with our cone
matrix it gives the negative of the triangle boundary. The
classification and addition statements agree after this uniform
sign change. See also \cite[Chapter 11]{p103}.
\end{proof}

\begin{proposition}[Classical Tor from a total left derived functor]
\label{res:R14}
Fix a ring \(R\) and a left \(R\)-module \(M\). For a right
\(R\)-module \(A\), \(H^{-n}\bigl(L(-\otimes_R M)(A[0])\bigr)\)
is naturally \(\operatorname{Tor}_n^R(A,M)\). In degree zero this is
\(A\otimes_R M\), and its connecting maps agree with the classical
ones with the chain indices reversed from cochain indices.
\end{proposition}
\begin{proof}
Choose a projective right-module chain resolution \(P_\bullet\to A\),
place \(P_n\) in degree \(-n\) with differential \(d_n:P_n\to P_{n-1}\),
and apply Theorem~\ref{res:R07}. On this model both constructions
take the homology of \(P_\bullet\otimes_R M\). Resolution comparison
makes the identification independent of the choice and natural in both
module variables. The balanced computation via a left-module resolution
is the comparison of Theorem~\ref{res:B05}. Right exactness gives
the degree-zero statement. For a short exact sequence of right modules
use a compatible horseshoe resolution. Its degreewise split sequence
remains exact after tensoring and gives both connecting maps by the
same lift and differential. The sign of the map in the derived
triangle is the one identified in Theorem~\ref{res:D04}. See
\cite[Section 3.2]{weibel} and \cite[Chapter 11]{p103}.
\end{proof}

\Needspace{24\baselineskip}
\section{Higher Yoneda extensions and products}

For objects \(X,Y\) of the abelian category and \(n\ge1\), an
\(n\)-extension of \(X\) by \(Y\) is an exact sequence
\[
  0\longrightarrow Y\longrightarrow E_{n-1}\longrightarrow
  \cdots\longrightarrow E_0\longrightarrow X\longrightarrow 0 .
\]
Take equivalence generated by maps of such exact sequences that
are the identity on both endpoints. The Baer operation is direct
sum followed by pullback at the right endpoint along the diagonal
of \(X\) and pushout at the left endpoint along the codiagonal of
\(Y\). For \(n=0\) use \(\operatorname{Hom}(X,Y)\). Splicing an
\(m\)-extension of \(Y\) by \(Z\) on the left of an \(n\)-extension
of \(X\) by \(Y\) uses the composite through \(Y\) at their junction.

\begin{theorem}[The higher Yoneda product]\label{res:R15}
In any abelian category the equivalence classes of \(n\)-extensions
of \(X\) by \(Y\), with the Baer operation, are naturally isomorphic
to \(\operatorname{Ext}^n(X,Y)\) as defined in
\ref{res:R11}. In degree one the isomorphism uses the positive
projection convention of \ref{res:R13}. Under these isomorphisms,
splicing a class \(\alpha\) of degree \(m\) on the left of a class
\(\beta\) of degree \(n\) is
\[
       \alpha[n]\circ\beta:X\longrightarrow Z[n+m],
       \qquad
       \beta:X\longrightarrow Y[n],\quad
       \alpha:Y\longrightarrow Z[m].
\]
Degree-zero classes act by pullback or pushout on extensions and by
ordinary composition in degree zero.
\end{theorem}

\begin{proof}
Put the terms from \(Y\) through \(E_0\) in degrees \(-n,\ldots,0\)
with the displayed differentials, forming a cochain complex \(C_E\).
The augmentation \(s_E:C_E\to X[0]\) is a quasi-isomorphism.
Projection \(p_E:C_E\to Y[n]\) is a chain map, equal to the identity
on \(Y\) in degree \(-n\). Write \(u_E=p_Es_E^{-1}\) for the
unscaled roof and put
\[
             y_E=(-1)^{n(n-1)/2}u_E .
\]
These expressions are invariant under maps of extensions with
identity endpoints, since the induced chain maps commute with
augmentation and projection and are quasi-isomorphisms.

We verify that the unscaled roofs already give a bijection of sets.
Represent a morphism \(X[0]\to Y[n]\) by a roof with middle complex
\(L\), using \ref{res:D02}. Apply the good truncations of
\ref{res:D05}, first \(\tau_{\le0}\), then \(\tau_{\ge -n}\).
They are quasi-isomorphisms in the relevant directions because
\(H^i(L)=0\) for \(i\ne0\). The arrow to \(Y[n]\) descends in
degree \(-n\) because it kills the incoming boundaries. Write
\(U\) for this truncated complex, \(s^0:U^0\to X\) for the
augmentation, and \(f:U^{-n}\to Y\) for the other arrow.
Then \(U^{-n}\to U^{-n+1}\) is monic, \(s^0\) is epic,
and the augmented row is exact. Push out the first monomorphism
along \(f\). Pushouts of monomorphisms in an abelian category
preserve the short exact row with the same cokernel, so the
resulting row is an \(n\)-extension by \(Y\). The natural map
from \(U\) to its complex is \(f\) at the left, the pushout map
next, and the identity in the remaining terms. It commutes with
both arrows of the roof and induces the identity on \(X\).
This proves surjectivity.

If two extension roofs agree, use the common refinement of
\ref{res:D02} and apply the same two truncations to it.
The maps to the extension complexes descend through the lower
truncation because their targets vanish below degree \(-n\).
The two maps to \(X[0]\) become identical: a homotopy has no contribution
in degree zero once the source is supported in degrees at most
zero. The two maps to \(Y[n]\) then differ in degree \(-n\) by a
map \(h d_U^{-n}\) for some \(h:U^{-n+1}\to Y\), by the chain
homotopy formula. The two pushouts along \(f\) and \(f+h d_U^{-n}\)
are isomorphic by the biproduct map
\((y,v)\mapsto(y-h(v),v)\). Indeed, it maps the defining relations
\((f(a),-d_U(a))\) to \((f(a)+h d_U(a),-d_U(a))\), and fixes
the endpoint maps. Each truncated refinement induces a map
from its pushout extension to the corresponding original
extension, fixing \(X\) and \(Y\). This gives the required
equivalence and proves injectivity.

Endpoint maps act on the roofs by precomposition and postcomposition:
pullback changes the augmentation by its pullback square, and
pushout changes the projection by its pushout square. The induced
maps of extension complexes give those two identities of roofs.
Taking direct sums and then the diagonal and codiagonal proves
additivity. Multiplying by the chosen sign preserves these
properties. In degree one the factor is \(1\), and \(p_Es_E^{-1}\)
is the positive projection of \ref{res:R13}, namely the negative
of the triangle boundary of \ref{res:D04}.

For the product, take \(n,m>0\). Let \(C\) be the complex of the
spliced extension of \(\alpha\) on the left of \(\beta\).
It maps to \(C_\beta\) by the identity on the \(\beta\)-terms,
by the augmentation of the rightmost \(\alpha\)-term onto \(Y\)
in degree \(-n\), and by zero to the left. Call this map \(a\).
It preserves the augmentation to \(X\). There is also a chain
map \(b:C\to C_\alpha[n]\) which is zero to the right of
degree \(-n\) and is multiplication by \((-1)^{nj}\) on the
\(\alpha\)-term in degree \(-n-j\) for \(0\le j\le m\).
The differential in \(C_\alpha[n]\) has the factor \((-1)^n\),
so these scalars commute with the differential. At degree \(-n\)
the scalar is \(1\), whence \(s_\alpha[n]b=p_\beta a\).
Since \(s_\beta a\) is the augmentation of the spliced extension,
the composite of the unscaled roofs is represented on \(C\) by
\(p_\alpha[n]b\). This equals \((-1)^{nm}\) times the positive
projection of \(C\) onto \(Z[n+m]\).

Finally,
\[
 \frac{(n+m)(n+m-1)}{2}
       =\frac{n(n-1)}{2}+\frac{m(m-1)}{2}+nm .
\]
It follows that the normalized classes give exactly shifted
composition. The cases with a degree-zero factor follow from
naturality. Thus the sign comparison uses the cochain shift
and the degree-one projection throughout.
For the extension classification see
\cite[Tags 06XP, 06XT and 06XU]{stacks-derived}.
The calculation above specifies the normalization for our
cochain shifts and extension differentials.
\end{proof}
